% Source: https://artofproblemsolving.com/community/c3h2029806_happy_pi_day_maybe_a_competition
% Stable identity: aops-topic-2029806-q01
Let the circle have radius \(r\). The diagonals of the inscribed square are diameters, of length \(2r\).

\begin{asy}
usepackage("amsmath");
size(3.5cm);
defaultpen(fontsize(10pt));
draw(circle((0,0),1));
pair A=dir(45),B=dir(135),C=dir(225),D=dir(315);
fill(A--B--C--D--cycle,paleblue); draw(A--B--C--D--cycle,linewidth(0.9));
draw((0,0)--A,red); label("$r$",0.5*A,SE,red); dot((0,0)); label("$K$",(0,0),S);
label("$A$",A,NE);label("$B$",B,NW);label("$C$",C,SW);label("$D$",D,SE);
label("\mbox{}",(min(currentpicture,user=true)+max(currentpicture,user=true))/2,black); // keeps the text after the figure black when LaTeX typesets the labels inline
\end{asy}

A square with diagonal \(2r\) has area \(\dfrac{(2r)^2}{2}=2r^2\). The circle has area \(\pi r^2\), so the fraction is
\[
\dfrac{2r^2}{\pi r^2}=\boxed{\dfrac{2}{\pi}}.
\]
